\chapter{IMPLEMENTATION DETAILS}


\section{Acoustic Threat Detection Implementation}

\subsection{Training Pipeline}
\label{subsec:training-pipeline}

The acoustic model development pipeline was implemented independently in
Python using TensorFlow/Keras, rather than the Edge Impulse platform
originally proposed. This gave full control over preprocessing, the
architecture search described in Section~\ref{subsec:model-evolution}, and
quantization -- control that proved necessary once memory profiling on
target hardware revealed how extensively the model had to be redesigned.

\textbf{Data Preparation:} The 5-class audio subset was split using a
stratified 80/20 train-test split per class (\texttt{split\_dataset.py}),
computed on the original, unaugmented files to keep the test set clean.
Augmentation (\texttt{augment\_training.py}) -- pitch shifting, time
stretching, and additive noise injection -- was applied only to the four
threat classes (axe, chainsaw, gunshot, handsaw) and deliberately withheld
from the background class, so that augmentation does not make background
clips sound more threat-like. This expanded the augmented training set to
roughly 360 files per threat class and 544 files for background.

\textbf{Feature Extraction:} Raw 16~kHz audio was converted to a log-Mel
spectrogram using a Hann window, 512-point FFT, and a Mel filterbank,
followed by per-frame normalisation (mean subtraction, unit standard
deviation). Figure~\ref{fig:example_spectrograms} shows example spectrograms
for representative classes, illustrating the harmonic bands of chainsaw
noise, the sharp broadband transient of a gunshot, and the impact profile of
an axe strike.

\begin{figure}[H]
    \centering
    \includegraphics[width=0.85\textwidth]{example_spectrograms.png}
    \caption{Log-Mel spectrograms for representative classes}
    \label{fig:example_spectrograms}
\end{figure}

\textbf{Model Training:} All three architecture variants (Section
\ref{subsec:model-evolution}) were trained using categorical cross-entropy
loss with the Adam optimizer. The final deployed (tiny) variant uses
\texttt{GlobalAveragePooling2D} in place of a flattened dense layer, avoiding
an oversized fully-connected weight matrix.

\subsection{INT8 Quantization}
\label{subsec:quantization-impl}

Each trained model was quantized to 8-bit integers using TensorFlow Lite's
post-training quantization API, with \texttt{TFLITE\_BUILTINS\_INT8} as the
target op set and INT8 fixed as both the inference input and output type:

\begin{verbatim}
converter = tf.lite.TFLiteConverter.from_keras_model(model)
converter.optimizations = [tf.lite.Optimize.DEFAULT]
converter.target_spec.supported_ops = \
    [tf.lite.OpsSet.TFLITE_BUILTINS_INT8]
converter.inference_input_type = tf.int8
converter.inference_output_type = tf.int8
converter.representative_dataset = representative_dataset
\end{verbatim}

A representative dataset of 100 randomly sampled training clips was used to
calibrate activation quantization ranges. For the final deployed tiny model
(retrained on the 5-class dataset), calibration produced an input scale of
0.076175 with zero-point $-13$ (mapping the INT8 range $[-128, 127]$ to
approximately $[-8.8, 10.7]$ in the normalised float domain), and an output
scale of 0.003906 with zero-point $-128$ for dequantizing softmax
probabilities. This reduced the tiny model to a 19.24~KB INT8 TFLite file
(\texttt{model\_int8.tflite}), compiled into a C header
(\texttt{model\_data.h}) and flashed to the ESP32-S3 firmware. The resulting
per-model file sizes and architecture parameters are summarised in
Table~\ref{tab:model-evolution-impl}.

\begin{table}[H]
    \centering
    \caption{Deployed model size and parameter count by architecture variant}
    \label{tab:model-evolution-impl}
    \begin{tabular}{lccc}
        \toprule
        \textbf{Metric} & \textbf{Full} & \textbf{Compact} & \textbf{Tiny (deployed)} \\
        \midrule
        Parameters          & $\sim$65{,}000 & $\sim$24{,}000 & $\sim$9{,}000 \\
        TFLite (INT8) size  & 78~KB          & 27~KB          & 19~KB \\
        Tensor arena        & (undeployable) & (undeployable) & 90~KB \\
        \bottomrule
    \end{tabular}
\end{table}



Table~\ref{tab:memory-budget} breaks down the full runtime memory budget on
the ESP32-S3. The 90~KB tensor arena and 19~KB of flash-resident model
weights leave sufficient headroom for the remaining audio pipeline,
FreeRTOS, and networking stack within the 512~KB SRAM budget.

\begin{table}[H]
    \centering
    \caption{ESP32-S3 runtime memory budget}
    \label{tab:memory-budget}
    \begin{tabular}{lcc}
        \toprule
        \textbf{Component} & \textbf{Size} & \textbf{Region} \\
        \midrule
        Model weights          & 19~KB & Flash \\
        Tensor arena           & 90~KB & DRAM \\
        Log-Mel buffer (int16) & 32~KB & Heap \\
        FFT + magnitude        & 5~KB  & Heap \\
        I2S DMA                & 8~KB  & Internal \\
        Ring buffer + stack    & $\sim$2~KB & Internal \\
        \bottomrule
    \end{tabular}
\end{table}

\section{LoRa Network Simulation Implementation using ns-3}

To evaluate the wireless communication capability of the proposed forest monitoring
system, a LoRaWAN-based simulation environment was developed using ns-3 with the
LDSE multi-hop routing protocol. The simulation environment leverages the ns-3.48
contrib modules:

\begin{itemize}
    \item \textbf{ns-3 lorawan module} (\texttt{ns-3.48/contrib/lorawan}/): Provides
          comprehensive LoRaWAN stack implementation including end devices, gateways,
          ALOHA-based contention, Adaptive Data Rate (ADR), Class A MAC operations,
          and application layer protocols.
    \item \textbf{ns-3 ldse module} (\texttt{ns-3.48/contrib/ldse}/): Implements the
          Layered Dynamic Synchronization Energy-saving protocol with core components:
          \begin{itemize}
              \item \texttt{ldse-sensor-app.cc/h}: Application layer for data generation
              \item \texttt{ldse-routing.cc/h}: Multi-hop routing with layer assignment
              \item \texttt{ldse-sync.cc/h}: Time synchronization via FTSP (Flooding Time
                Synchronisation Protocol)
              \item \texttt{ldse-layering.cc/h}: Hierarchical layer topology formation
              \item \texttt{ldse-channel-manager.cc/h}: Multi-channel separation (public/private)
              \item \texttt{ldse-energy-controller.cc/h}: Energy consumption modeling
              \item \texttt{ldse-epoch-scheduler.cc/h}: Epoch structure management
          \end{itemize}
    \item \textbf{Channel Model:} Realistic forest path loss model combining
          log-distance path loss with foliage attenuation factors to simulate how
          thick trees degrade LoRa signals compared to open field conditions.
\end{itemize}

The simulation environment consists of:

\begin{itemize}
    \item LoRa end devices configured with LDSE multi-hop routing.
    \item LoRa gateway node at the center of the network area.
    \item LDSE protocol stack managing layer assignment, time synchronization, and
          relay forwarding.
    \item Forest channel model with foliage attenuation and log-distance path loss.
    \item Network communication components provided by the INET framework.
\end{itemize}

A basic LoRa communication scenario was implemented where LoRa sensor nodes
communicate through the LDSE multi-hop routing protocol to the gateway.

The simulation workflow includes:

\begin{enumerate}
    \item Creation and configuration of LoRa nodes using the ns-3 lorawan helper.
    \item Configuration of LDSE protocol parameters (layer assignment, time slots,
          guard intervals, parent selection).
    \item Configuration of forest channel model (foliage attenuation, log-distance
          path loss exponent).
    \item Packet generation from end devices (sensor data + alert payloads).
    \item Wireless packet transmission through the LoRa medium with forest propagation.
    \item LDSE multi-hop forwarding: intermediate relay nodes receive and forward
          packets based on layer assignment and parent selection logic.
    \item Reception and processing of packets at the gateway.
    \item Performance metric extraction: PDR, end-to-end latency, energy consumption,
          duty cycle compliance, and collision rate.
\end{enumerate}

The implemented LoRa network topology in ns-3 is illustrated in
Fig.~\ref{fig:lora_simulation_ns3}.

\begin{figure}[!ht]
    \centering
    \includegraphics[width=0.75\textwidth]{lora_simulation.png}
    \caption{LoRaWAN communication using ns-3 with LDSE multi-hop routing}
    \label{fig:lora_simulation_ns3}
\end{figure}


\subsection{Basic LoRa Node Communication Implementation}

Before implementing the proposed multi-hop routing mechanism, a basic LoRa
communication scenario was developed to verify the operation of individual LoRa
end devices. In this implementation, two LoRa sensor nodes were configured to
exchange data through the LoRa wireless communication medium.

The first node acts as a LoRa transmitter, generating application data packets
at predefined time intervals. These packets are processed through the LoRa
protocol stack and transmitted using the LoRa physical layer. The second node
acts as a LoRa receiver, which receives the transmitted packets and processes
the incoming data.

The communication process consists of the following steps:

\begin{enumerate}
    \item The LoRa transmitter node generates sensing data packets.
    \item The packets are passed through the application and LoRa protocol layers.
    \item The LoRa radio module performs wireless transmission using configured
    parameters such as spreading factor, bandwidth, and transmission power.
    \item The receiving LoRa node detects and receives the transmitted packet.
    \item The received data is delivered to the application layer for further
    processing.
\end{enumerate}

This basic communication model validates the successful exchange of information
between LoRa devices and confirms the proper configuration of the FLORA
framework and INET communication modules. The implemented two-node LoRa
communication scenario serves as the initial building block for extending the
network towards gateway-based and multi-hop communication.

The implemented two-node LoRa communication topology is shown in
Fig.~\ref{fig:two_node_lora}.

\begin{figure}[!ht]
    \centering
    \includegraphics[width=0.75\textwidth]{two_node_lora.png}
    \caption{LoRa communication between two LoRa end devices in OMNeT++}
    \label{fig:two_node_lora}
\end{figure}